% ============================================================================
% AI-DRAFTED BLOCK (Claude), revision 3, 2026-08-23. Entire file is one
% generated block, delimited per the project seam protocol. Author passages
% pending:
%   (1) the economic-mechanism paragraph in Section 1 is marked for redraft
%       by the author per the project drafting protocol;
%   (2) bibliography entries marked % TODO carry unverified volume/page data.
% Revision 3 changes vs revision 2: simulation-backed upgrades from the
% companion study (through tag tgarch-study-round3-v1). Remark on filtering
% now carries the two closed-form error laws with measured constants; the
% tails remark states the verified Kesten-to-inverse-gamma index
% convergence instead of the former non-commuting-limits claim; Section 5
% reports the completed verification numbers, the identification results,
% and the fixed-budget shortfall illustration of the open uniform
% integrability problem. The full-budget clean-tag pricing rerun and the
% nested path-budget diagnostic are complete.
% Revision 2 changes vs revision 1: notation aligned with the published
% log-normal beta SV paper (signed volatility beta, residual vol-of-vol,
% hatted Q-parameters); Student-t section removed; one-step kernels presented
% as a product measure change; exact-filtration proposition added; remarks
% added on the minimal kernel, on the Girsanov (not Ito) origin of the
% quadratic term, and on the correspondence with the paper's parameter
% transformation.
% ============================================================================
\documentclass[11pt]{article}
\usepackage[margin=1.1in]{geometry}
\usepackage{amsmath,amssymb,amsthm}
\usepackage{booktabs}

\newtheorem{theorem}{Theorem}
\newtheorem{proposition}[theorem]{Proposition}
\newtheorem{lemma}[theorem]{Lemma}
\newtheorem{corollary}[theorem]{Corollary}
\newtheorem{assumption}[theorem]{Assumption}
\theoremstyle{remark}
\newtheorem{remark}[theorem]{Remark}

\newcommand{\E}{\mathbb{E}}
\newcommand{\Pm}{\mathbb{P}}
\newcommand{\Qm}{\mathbb{Q}}
\newcommand{\R}{\mathbb{R}}
\newcommand{\Dt}{\Delta}

\title{A Discrete-Time Origin of the Log-Normal Beta Stochastic Volatility
Model with Quadratic Drift}
\author{Working note (draft)\thanks{AI-assisted draft, revision 3, prepared
2026-08-23. All numbered claims carry proofs in the text. Bibliography
entries marked in the source require verification before any submission
use.}}
\date{August 23, 2026}

\begin{document}
\maketitle

\begin{abstract}
The log-normal beta stochastic volatility model with quadratic drift of Sepp
and Rakhmonov (2023) connects its statistical and risk-neutral dynamics
through a Girsanov change with two risk prices proportional to volatility,
whose constants $\bar\lambda_0$ and $\bar\lambda_1$ are free parameters.
A discrete-time model generating these dynamics, and an economic
identification of the two constants, have been missing.
We apply one-period exponential-quadratic pricing kernels with
volatility-dependent loadings to a threshold recursion on the volatility
level driven by the single return innovation.
The risk-neutral drift reproduces the parameter transformation of the
continuous model with $\bar\lambda_0=\gamma_1$, the slope of the Sharpe
ratio in volatility, and $\bar\lambda_1=-(m_1/s_1)\,\eta_1$, the slope of
the kernel variance preference.
The drift intercept $\kappa_1\theta$ is invariant under every admissible
kernel, and the recursion converges weakly to the continuous model under
both measures.
The recursion supplies an exact volatility filter from returns at every
step and a cross-measure restriction for joint estimation from returns and
options.
\end{abstract}

\section{Introduction and results}\label{sec:intro}

The log-normal beta stochastic volatility (SV) model with quadratic drift
specifies the volatility dynamics
$d\sigma_t=(\kappa_1+\kappa_2\sigma_t)(\theta-\sigma_t)\,dt
+\beta\sigma_t\,dW^{(0)}_t+\varepsilon\sigma_t\,dW^{(1)}_t$, with the
volatility beta $\beta\in\R$ measuring the response of volatility to
returns and with $\varepsilon>0$ the residual volatility of volatility
\cite{SeppRakhmonov2023}.
The published paper connects the statistical and the risk-neutral dynamics
through a Girsanov change with risk prices
$\lambda_0(t)=\bar\lambda_0\sigma_t$ and
$\lambda_1(t)=\bar\lambda_1\sigma_t$, which shifts the quadratic
mean-reversion rate to
$\hat\kappa_2=\kappa_2+\beta\bar\lambda_0+\varepsilon\bar\lambda_1$ and
leaves the products $\kappa_1\theta=\hat\kappa_1\hat\theta$ matched.
Carr and Willems (2019) study a close relative with the same drift
structure.
The discrete-time side of this model family has remained open.
Nelson (1990) showed that the standard GARCH(1,1) recursion converges to a
variance diffusion with linear mean reversion, not to the Heston model.
Duan (1997) extended the convergence theory to an augmented family that
nests recursions on the volatility level, whose limits carry linear
drift.
The affine recursion of Heston and Nandi (2000) converges to the Heston
model with perfectly correlated shocks.
No discrete-time model has been shown to generate the quadratic drift, and
the continuous-time constants $\bar\lambda_0$ and $\bar\lambda_1$ have had
no primitive interpretation.

We establish three results.
First, we derive the transformation of the volatility drift under
one-period exponential-quadratic pricing kernels with volatility-dependent
loadings (Theorem~\ref{thm:main}).
The physical drift polynomial $d_0+d_1\sigma+d_2\sigma^2$ maps to
$d_0+(d_1+\Lambda_0)\sigma+(d_2+\Lambda_1)\sigma^2$, where $\Lambda_0$ and
$\Lambda_1$ collect the intercepts and the slopes of the Sharpe ratio
$\gamma(\sigma)=\gamma_0+\gamma_1\sigma$ and of the kernel variance
preference $\eta(\sigma)=\eta_0+\eta_1\sigma$.
The map reproduces the parameter transformation of the continuous model
exactly, with the induced risk prices $\lambda_0(\sigma)=\gamma(\sigma)$
and $\lambda_1(\sigma)=-(m_1/s_1)\,\eta(\sigma)$.
The free constants of the published paper acquire discrete primitives:
$\bar\lambda_0$ is the Sharpe-ratio slope and $\bar\lambda_1$ is the slope
of the variance preference scaled by the news-impact constant $m_1/s_1$
(Corollary~\ref{cor:theta} and Remark~\ref{rem:match}).
Second, we show that the drift intercept $d_0=\kappa_1\theta$ is invariant
under every admissible pricing kernel, at every step size and in the limit
(Proposition~\ref{prop:invariance}).
The invariance turns the matching convention
$\kappa_1\theta=\hat\kappa_1\hat\theta$ of the published paper into a
structural restriction available for joint estimation.
Third, we prove weak convergence of the recursion to the continuous model
under both measures (Theorem~\ref{thm:limit}), and we show that the
recursion admits an exact volatility filter from returns at every step
(Proposition~\ref{prop:filter}), a property the two-Brownian limit does not
share.

% [AUTHOR: mechanism paragraph below is a placeholder draft. Per project
% protocol the economic-mechanism passage is author-written. Redraft.]
Under the risk-neutral measure the volatility drift acquires the covariance
between the pricing kernel and the volatility innovation.
The kernel prices two features of the single return innovation: its level,
at the market Sharpe ratio, and its size, at the variance preference.
When the Sharpe ratio rises with volatility, the compensation for the level
shock scales with $\sigma$, and its transmission into the volatility
recursion scales with $\beta\sigma$, which contributes
$\beta\gamma_1\sigma^2$ to the drift.
The variance-preference slope acts in the same way through the size channel
carried by $\varepsilon$.
The quadratic drift is the continuous-time image of level-dependent risk
prices, and the sign of $\beta$ decides which channel can fund it.

The construction subsumes the linear-drift case: setting
$\kappa_2=\gamma_1=\eta_1=0$ recovers a linear drift under both measures,
and the physical recursion is exactly the threshold model of Zakoian (1994)
written on the volatility level.
As a discrete-to-continuous bridge, the construction substitutes for Heston
and Nandi (2000) on assets with log-normal volatility and a signed
volatility beta.
The cost of switching is the loss of an exact discrete moment-generating
function recursion, and the retained benefit is a single-shock recursion
whose volatility is an exact function of observed returns.

One problem stays open, and we state it rather than hide it.
Weak convergence delivers option prices only for bounded payoffs, and call
prices require uniform integrability of the discrete price family, which
fails for $\hat\kappa_2=0$ with $\beta>0$ because the discounted limiting
stock-price process is a strict local martingale.
Section~\ref{sec:open} formulates the required uniform moment bound, which
we regard as the central technical claim of any full paper built on this
note.
Section~\ref{sec:model} defines the discrete model.
Section~\ref{sec:kernel} derives the risk-neutral recursion.
Section~\ref{sec:limit} states the diffusion limits, and
Section~\ref{sec:open} collects verification steps, constraints, and the
open problem.

\section{The discrete model under the physical measure}\label{sec:model}

We fix a horizon $T$, a step $\Dt>0$, and the grid $t_k=k\Dt$.
Innovations $z_{k}$ are independent standard Gaussian variables on a
filtered space with filtration $\mathcal{F}_k$ generated by the
innovations.
We write $m_1=\E|z|=\sqrt{2/\pi}$ and
$s_1^2=\operatorname{Var}(|z|)=1-2/\pi$, and we define the standardized
size innovation
\begin{equation}\label{eq:w}
w \;=\; \frac{|z|-m_1}{s_1},
\end{equation}
which satisfies $\E[w]=0$, $\E[w^2]=1$ and $\E[z\,w]=0$ by symmetry of $z$.
The physical model for the log price $X=\ln S$ and the volatility $\sigma$
is
\begin{align}
X_{k+1}&=X_k+\Big(r+\gamma(\sigma_k)\,\sigma_k
-\tfrac{1}{2}\sigma_k^2\Big)\Dt+\sigma_k\sqrt{\Dt}\,z_{k+1},
\label{eq:Xdyn}\\
\sigma_{k+1}&=\sigma_k+b(\sigma_k)\,\Dt
+\sigma_k\sqrt{\Dt}\,\big(\beta\,z_{k+1}+\varepsilon\,w_{k+1}\big),
\label{eq:sigdyn}
\end{align}
with the quadratic drift and the Sharpe ratio
\begin{equation}\label{eq:drift}
b(\sigma)=(\kappa_1+\kappa_2\sigma)(\theta-\sigma)
=d_0+d_1\sigma+d_2\sigma^2,
\qquad
\gamma(\sigma)=\gamma_0+\gamma_1\sigma ,
\end{equation}
so that $d_0=\kappa_1\theta>0$, $d_1=\kappa_2\theta-\kappa_1$ and
$d_2=-\kappa_2\le0$, with $\kappa_1,\theta>0$ and $\kappa_2\ge0$ as in the
continuous model.
The volatility beta $\beta\in\R$ is signed, the residual volatility of
volatility satisfies $\varepsilon>0$, and the total variance of the
volatility innovation is $\vartheta^2=\beta^2+\varepsilon^2$, matching the
continuous convention.
The excess return in \eqref{eq:Xdyn} equals $\gamma(\sigma)\sigma$, so
$\gamma(\sigma)$ is the conditional Sharpe ratio of the asset.
We believe the affine Sharpe specification is the right trade-off because
it is the minimal specification that can move the quadratic drift
coefficient and it nests the proportional risk price
$\lambda_0=\bar\lambda_0\sigma$ of the continuous model at $\gamma_0=0$
with $\bar\lambda_0=\gamma_1$.

\begin{lemma}\label{lem:innov}
The volatility innovation admits the threshold form
$\beta z+\varepsilon w
=(\beta+\varepsilon/s_1)\,z^{+}-(\beta-\varepsilon/s_1)\,z^{-}
-\varepsilon m_1/s_1$
with $z^{+}=\max(z,0)$ and $z^{-}=\max(-z,0)$, and its variance equals
$\vartheta^2$.
\end{lemma}

\begin{proof}
Substitute $z=z^{+}-z^{-}$ and $|z|=z^{+}+z^{-}$ into
$\beta z+\varepsilon w$.
The variance follows from $\E[zw]=0$.
\end{proof}

The threshold form identifies \eqref{eq:sigdyn} as a Zakoian (1994) model
on the volatility level with an intercept, so the physical dynamics require
no new econometrics, and for $\beta<0$ the down-move slope
$\varepsilon/s_1-\beta$ exceeds the up-move slope, which is the leverage
asymmetry.
Both equations are driven by functions of the single innovation $z$, and
the next proposition records the consequence that the volatility path is an
exact function of the return path.

\begin{proposition}\label{prop:filter}
Fix the model parameters and $\sigma_0>0$, and suppose the recursion is
modified by the localization device of Remark~\ref{rem:trunc} so that
$\sigma_k>0$ for all $k$.
Then for every $k$ the volatility $\sigma_k$ is measurable with respect to
$\sigma(X_1,\dots,X_k;\sigma_0)$: the innovation is recovered from the
observed return by
\begin{equation}\label{eq:filter}
z_{k+1}
=\frac{X_{k+1}-X_k-\big(r+\gamma(\sigma_k)\sigma_k
-\tfrac{1}{2}\sigma_k^2\big)\Dt}{\sigma_k\sqrt{\Dt}},
\end{equation}
and \eqref{eq:sigdyn} then determines $\sigma_{k+1}$.
The same statement holds under the risk-neutral measure with the
risk-neutral drifts.
\end{proposition}

\begin{proof}
Induction over $k$: given $\sigma_k>0$, equation \eqref{eq:filter} inverts
\eqref{eq:Xdyn} for $z_{k+1}$, the size innovation $w_{k+1}$ is the
deterministic transform \eqref{eq:w} of $z_{k+1}$, and \eqref{eq:sigdyn}
delivers $\sigma_{k+1}$.
\end{proof}

\begin{remark}[Estimation]\label{rem:qmle}
Proposition~\ref{prop:filter} makes estimation identical to threshold GARCH
inference.
The volatility $\sigma_k$ is predictable, so the return over
$(t_k,t_{k+1}]$ is conditionally Gaussian with mean
$(r+\gamma(\sigma_k)\sigma_k-\tfrac{1}{2}\sigma_k^2)\Dt$ and variance
$\sigma_k^2\Dt$, and the sample likelihood is the prediction-error
decomposition
\begin{equation}\label{eq:qmle}
\ell=-\frac{1}{2}\sum_{k=0}^{n-1}
\Big(\ln\big(2\pi\sigma_k^2\Dt\big)+z_{k+1}^2\Big),
\end{equation}
with $z_{k+1}$ recovered by \eqref{eq:filter} and $\sigma_{k+1}$ updated by
\eqref{eq:sigdyn} inside the same loop.
The Sharpe function sits in the conditional mean, so $\gamma_1$ is
estimated as an in-mean coefficient from returns alone, and it is the same
constant that the kernel maps to the risk price $\bar\lambda_0$ in
Section~\ref{sec:kernel}.
The initial value $\sigma_0$ is a parameter or is set to $\theta$, with an
influence that decays at the mean-reversion rate.
Under Gaussian innovations \eqref{eq:qmle} is the exact likelihood, and
otherwise it is a quasi-likelihood in the usual sense.
\end{remark}

\begin{remark}\label{rem:filtergap}
The limit model does not share this property, because the size innovations
$w_k$ aggregate into a Brownian motion $W^{(1)}$ independent of the return
Brownian $W^{(0)}$, and the limit volatility is not spanned by returns.
Exact filterability at every $\Dt$ and unspanned volatility in the limit
are reconciled by the asymptotic filtering theory of Nelson and Foster
(1994), and the reconciliation is quantitative.
Applied to a path of the two-shock limit model, the filter of
Proposition~\ref{prop:filter} corrects its own error at the data-driven
rate
\begin{equation}\label{eq:filterrate}
g_\Dt=\kappa_{\mathrm{lin}}
+\frac{\varepsilon\,m_1}{s_1}\,\frac{1}{\sqrt{\Dt}},
\qquad
\kappa_{\mathrm{lin}}=\kappa_1+\kappa_2\theta,
\end{equation}
because an overestimated volatility shrinks the recovered residuals and
the size channel pushes the estimate back down.
Balancing this gain against the unspanned noise gives the steady-state
error
\begin{equation}\label{eq:filterrmse}
\mathrm{RMSE}\;\approx\;\bar\sigma\,
\sqrt{\varepsilon\, s_1/m_1}\;\Dt^{1/4},
\end{equation}
with $\bar\sigma$ the average volatility level, which is the
Nelson--Foster rate with explicit constants.
The companion simulation study confirms both
formulas.\footnote{All simulation numbers quoted in this note are from the
companion study (module \texttt{volatility\_book.ch\_discrete\_vol}). The
full-budget clean rerun and closing diagnostics use tag
\texttt{tgarch-study-round3-v1}; earlier inputs retain their producing tags.
The file \texttt{cited\_numbers.json} maps every quoted value to its JSON
pointer and producing tag, with seeds, package versions and hashes retained
in the study archive.}
Fitted error exponents are $0.249$ and $0.234$ at the two parameter sets
against the predicted $1/4$, the gain slope matches
$\varepsilon m_1/s_1$ within a tenth of a percent, and the level matches
\eqref{eq:filterrmse} within ten percent using the realized average
volatility.
A pathwise log-decay regression measures the intercept as the
Jensen-corrected value
$\kappa_{\mathrm{lin}}+\tfrac{1}{2}(\varepsilon m_1/s_1)^2$, within about nine
percent, because the multiplicative noise of the error accelerates the
decay of its logarithm.
\end{remark}

\begin{table}[h]
\centering
\caption{Notation. Quantities carry units of inverse time where the
continuous model requires it ($r$, $\kappa_1$, $\kappa_2\theta$, $d_i$),
and $\Dt$ carries units of time.}
\begin{tabular}{ll}
\toprule
$z,\,w$ & return innovation, size innovation \eqref{eq:w}\\
$m_1,\,s_1$ & mean and standard deviation of $|z|$\\
$\beta,\,\varepsilon,\,\vartheta$ & volatility beta ($\beta\in\R$),
residual vol-of-vol, $\vartheta^2=\beta^2+\varepsilon^2$\\
$b,\,\hat b$ & drift polynomials of $\sigma$ under $\Pm$ and $\Qm$\\
$d_0,d_1,d_2$ and $\hat d_0,\hat d_1,\hat d_2$ & coefficients of the drift
polynomials\\
$\gamma(\sigma)=\gamma_0+\gamma_1\sigma$ & conditional Sharpe ratio\\
$\eta(\sigma)=\eta_0+\eta_1\sigma$ & variance-preference loading\\
$\Lambda(\sigma)=\Lambda_0+\Lambda_1\sigma$ & volatility risk adjustment
(Proposition~\ref{prop:qmoments})\\
$\kappa_1,\kappa_2,\theta$ and $\hat\kappa_1,\hat\kappa_2,\hat\theta$ &
drift parameters under $\Pm$ and $\Qm$\\
$\kappa_{\mathrm{lin}}=\kappa_1+\kappa_2\theta$ & linearized reversion
speed \eqref{eq:filterrate}\\
$\lambda_0(\sigma),\,\lambda_1(\sigma)$ & induced risk prices of the return
and size channels\\
\bottomrule
\end{tabular}
\end{table}

\section{Pricing kernels and the risk-neutral recursion}\label{sec:kernel}

This section performs the change of measure and computes the risk-neutral
drift of the volatility recursion.
The measure change is specified through one-period kernels, and we first
make precise in what sense this is a single change of measure rather than a
new one at every step.
Let $M_{k+1}$ be a positive $\mathcal{F}_{k+1}$-measurable variable with
$\E_k[M_{k+1}]=1$, and define
\begin{equation}\label{eq:product}
\zeta_n=\prod_{k=0}^{n-1}M_{k+1},
\qquad
\frac{d\Qm}{d\Pm}\Big|_{\mathcal{F}_n}=\zeta_n .
\end{equation}
The process $\zeta$ is a positive $\Pm$-martingale by construction, the
one-period kernels are its multiplicative increments, and \eqref{eq:product}
is the discrete counterpart of the It\^{o} exponential that defines the
measure change of the continuous model.
The factorization into one-period kernels is not an extra assumption,
because every positive martingale density factorizes this way.
The one-period form is what makes the discounted price a $\Qm$-martingale
over each step rather than only at the horizon, which is the discrete
statement of dynamic no-arbitrage.
The kernel class is exponential-quadratic in the return innovation, in the
spirit of the variance-dependent kernel of Christoffersen, Heston and
Jacobs (2013).
Its loadings are affine in the volatility level,
\begin{equation}\label{eq:kernel}
M_{k+1}
\;\propto\;
\exp\!\big(a_k z_{k+1}+b_k z_{k+1}^2\big),
\qquad
b_k=\sqrt{\Dt}\,\eta(\sigma_k),\quad
\eta(\sigma)=\eta_0+\eta_1\sigma ,
\end{equation}
normalized to unit conditional expectation, with $a_k$ determined below by
the martingale condition and with $b_k<1/2$, which holds for all small
$\Dt$ on compact volatility sets.
The $\sqrt{\Dt}$ scaling of the loadings is a choice, and any convergence
statement below is a statement about this scaling scheme, in line with the
non-uniqueness results of Corradi (2000).
We adopt it because it is the unique power scaling under which the risk
adjustment neither vanishes nor explodes in the limit.

\begin{lemma}\label{lem:qlaw}
Under $\Qm$ the innovation $z_{k+1}$ given $\mathcal{F}_k$ is Gaussian with
mean $\lambda_k=a_k\varsigma_k^2$ and variance
$\varsigma_k^{2}=(1-2b_k)^{-1}$.
\end{lemma}

\begin{proof}
The conditional $\Qm$-density is proportional to
$\varphi(z)\exp(a_k z+b_k z^2)
\propto\exp\{-\tfrac{1}{2}(1-2b_k)z^2+a_k z\}$, and completing the square
gives the stated Gaussian law.
\end{proof}

\begin{proposition}\label{prop:pinned}
For each $b_k<1/2$ there exists a unique $a_k$ such that the discounted
price $e^{-r t_k+X_k}$ is a $\Qm$-martingale, and it satisfies
\begin{equation}\label{eq:lambda}
\lambda_k=-\sqrt{\Dt}\,\gamma(\sigma_k)
-\tfrac{1}{2}\,\sigma_k\sqrt{\Dt}\,\big(\varsigma_k^{2}-1\big)
=-\sqrt{\Dt}\,\gamma(\sigma_k)+O(\Dt),
\end{equation}
locally uniformly in $\sigma_k$.
\end{proposition}

\begin{proof}
By Lemma~\ref{lem:qlaw} the martingale condition
$\E^{\Qm}_k[e^{X_{k+1}-X_k}]=e^{r\Dt}$ reads
$\gamma(\sigma)\sigma\Dt+\sigma\sqrt{\Dt}\,\lambda_k
+\tfrac{1}{2}\sigma^2\Dt(\varsigma_k^2-1)=0$, which solves to
\eqref{eq:lambda}.
The map $a\mapsto\lambda=a\varsigma_k^2$ is a strictly increasing bijection
of $\R$ for fixed $b_k<1/2$, so $a_k$ exists and is unique.
The remainder in \eqref{eq:lambda} is $O(\Dt)$ because
$\varsigma_k^2-1=2b_k/(1-2b_k)=2\sqrt{\Dt}\,\eta(\sigma)+O(\Dt)$, with
constants uniform on compact volatility sets since $\gamma$ and $\eta$ are
affine.
\end{proof}

The martingale condition therefore pins the level loading to the Sharpe
ratio, and the only free input of the kernel is the variance preference
$\eta$.
The proof also records a cancellation worth naming.
The It\^{o} compensator $-\sigma^2/2$ of the log return cancels against the
quadratic term of the risk-neutral moment generating function up to order
$\Dt^{3/2}$, so the convexity of log returns leaves no trace in
\eqref{eq:lambda} at the relevant order.
The next step is the one place in the derivation where something can go
wrong, so we spell out why it does not.
The volatility innovation contains $|z|$, and a location shift of size
$\sqrt{\Dt}$ in $z$ could in principle feed the Sharpe ratio into the size
channel at order $\sqrt{\Dt}$, which would entangle the two risk prices in
the drift.
It cannot, because the first derivative of the absolute moment with respect
to a location shift vanishes at zero for every symmetric law: with
$h(x)=\E|z+x|$ one has $h'(0)=\E[\operatorname{sgn}(z)]=0$.
Symmetry of the innovation law is load-bearing here, and the channels
decouple at leading order precisely because of it.
Under skewed innovations the decomposition below acquires a cross term
proportional to $h'(0)$, which we leave to a skewed extension.

\begin{proposition}\label{prop:qmoments}
Let $h(x)=2\varphi(x)+x\,(2\Phi(x)-1)$ denote the absolute moment of a unit
Gaussian shifted by $x$, with $\varphi$ and $\Phi$ the standard normal
density and distribution function.
Under the kernel \eqref{eq:kernel} with $a_k$ from
Proposition~\ref{prop:pinned},
\begin{equation}\label{eq:Lambda}
\E^{\Qm}_k\big[\beta z_{k+1}+\varepsilon w_{k+1}\big]
=\sqrt{\Dt}\,\Lambda(\sigma_k)+O(\Dt),
\qquad
\Lambda(\sigma)
=-\beta\,\gamma(\sigma)+\varepsilon\,\frac{m_1}{s_1}\,\eta(\sigma),
\end{equation}
together with
$\operatorname{Var}^{\Qm}_k(\beta z+\varepsilon w)
=\vartheta^2+O(\sqrt{\Dt})$ and
$\operatorname{Cov}^{\Qm}_k(z,\beta z+\varepsilon w)
=\beta+O(\sqrt{\Dt})$,
all locally uniformly in $\sigma_k$.
\end{proposition}

\begin{proof}
By Lemma~\ref{lem:qlaw},
$\E^{\Qm}_k|z|=\varsigma_k\,h(\lambda_k/\varsigma_k)$.
Direct differentiation gives $h'(x)=2\Phi(x)-1$ and $h''(x)=2\varphi(x)$,
so $h(x)=m_1+\varphi(0)x^2+O(x^4)$ and
$\E^{\Qm}_k|z|=\varsigma_k m_1+O(\lambda_k^2)
= m_1+\sqrt{\Dt}\,\eta(\sigma_k)\,m_1+O(\Dt)$.
Combining with $\E^{\Qm}_k[z]=\lambda_k$ from \eqref{eq:lambda} inside
\eqref{eq:w} yields \eqref{eq:Lambda}.
The variance and covariance are smooth functions of
$(\lambda_k,\varsigma_k)$ that equal $(\vartheta^2,\beta)$ at $(0,1)$, and
$(\lambda_k,\varsigma_k)=(0,1)+O(\sqrt{\Dt})$, which gives the stated
perturbations.
Uniformity holds on compact volatility sets because the loadings are
affine.
\end{proof}

\begin{theorem}\label{thm:main}
Under the kernel \eqref{eq:kernel} the volatility recursion under $\Qm$
takes the form
\begin{equation}\label{eq:qrec}
\sigma_{k+1}=\sigma_k+\hat b_{\Dt}(\sigma_k)\,\Dt
+\sigma_k\sqrt{\Dt}\,\bar{u}_{k+1},
\qquad
\bar{u}_{k+1}=\beta z_{k+1}+\varepsilon w_{k+1}
-\E^{\Qm}_k\big[\beta z_{k+1}+\varepsilon w_{k+1}\big],
\end{equation}
where $\bar u$ has conditional mean zero, conditional variance
$\vartheta^2+O(\sqrt{\Dt})$, conditional covariance $\beta+O(\sqrt{\Dt})$
with the return innovation, and
\begin{equation}\label{eq:driftmap}
\hat b_{\Dt}(\sigma)\;\longrightarrow\;
\hat b(\sigma)= d_0+\big(d_1+\Lambda_0\big)\sigma
+\big(d_2+\Lambda_1\big)\sigma^{2},
\qquad
\begin{aligned}
\Lambda_0&=-\beta\gamma_0+\varepsilon\, m_1\eta_0/s_1,\\
\Lambda_1&=-\beta\gamma_1+\varepsilon\, m_1\eta_1/s_1,
\end{aligned}
\end{equation}
locally uniformly as $\Dt\to0$.
In particular, the risk-neutral quadratic mean-reversion rate is
\begin{equation}\label{eq:kappa2}
\hat\kappa_2 \;=\; -\hat d_2
\;=\; \kappa_2+\beta\,\gamma_1-\varepsilon\,\frac{m_1}{s_1}\,\eta_1 .
\end{equation}
\end{theorem}

\begin{proof}
Adding and subtracting the conditional mean in \eqref{eq:sigdyn} gives
\eqref{eq:qrec} with the exact identity
$\hat b_{\Dt}(\sigma)=b(\sigma)
+\sigma\,\E^{\Qm}_k[\beta z+\varepsilon w]/\sqrt{\Dt}$, and
Proposition~\ref{prop:qmoments} gives
$\E^{\Qm}_k[\beta z+\varepsilon w]
=\sqrt{\Dt}\Lambda(\sigma)+O(\Dt)$ locally uniformly.
Collecting powers of $\sigma$ in $b(\sigma)+\sigma\Lambda(\sigma)$ yields
\eqref{eq:driftmap}, and \eqref{eq:kappa2} follows from $d_2=-\kappa_2$.
\end{proof}

\begin{remark}\label{rem:match}
Theorem~\ref{thm:main} lands exactly on the continuous transformation of
Sepp and Rakhmonov (2023, Section 2).
Their Girsanov change with risk prices $\lambda_0(\sigma)$ and
$\lambda_1(\sigma)$ shifts the volatility drift by
$-\sigma\,(\beta\lambda_0(\sigma)+\varepsilon\lambda_1(\sigma))$, and
comparison with \eqref{eq:Lambda} identifies the induced prices
\begin{equation}\label{eq:prices}
\lambda_0(\sigma)=\gamma(\sigma),
\qquad
\lambda_1(\sigma)=-\frac{m_1}{s_1}\,\eta(\sigma).
\end{equation}
Their proportional specification
$\lambda_0=\bar\lambda_0\sigma$, $\lambda_1=\bar\lambda_1\sigma$
corresponds to $\gamma_0=\eta_0=0$ with $\bar\lambda_0=\gamma_1$ and
$\bar\lambda_1=-(m_1/s_1)\,\eta_1$, and then \eqref{eq:kappa2} reads
$\hat\kappa_2=\kappa_2+\beta\bar\lambda_0+\varepsilon\bar\lambda_1$, which
is their transformation verbatim.
The constants gain discrete primitives: the equity risk price is the
Sharpe-ratio slope, the volatility risk price is the variance-preference
slope, and the conversion factor $m_1/s_1=\sqrt{2/\pi}/\sqrt{1-2/\pi}$ is a
property of the absolute-value news-impact function.
The intercepts $\gamma_0$ and $\eta_0$ move only the linear coefficient
$\hat d_1$, which explains why the proportional specification of the
published paper leaves $\kappa_1-\kappa_2\theta$ unchanged.
\end{remark}

\begin{remark}\label{rem:esscher}
The variance loading is optional.
With $\eta\equiv0$ the kernel is linear in the innovation, the minimal
one-period martingale change, and the limit exists unchanged with
$\bar\lambda_1=0$ and $\hat\kappa_2=\kappa_2+\beta\gamma_1$.
The quadratic loading exists for one purpose only, which is to generate a
nonzero price of the size channel, and it is not needed for the log-return
limit.
\end{remark}

\begin{remark}\label{rem:ito}
The quadratic term in \eqref{eq:kappa2} is a Girsanov product, not an
It\^{o} correction.
Setting $\gamma_1=\eta_1=0$ gives $\hat\kappa_2=\kappa_2$ exactly, although
every convexity term of the model, including the log-return compensator
$-\sigma^2/2$, remains in place, and the proof of
Proposition~\ref{prop:pinned} shows that this compensator cancels at the
relevant order.
The shift $\hat\kappa_2-\kappa_2$ is the product of the level-proportional
dispersion $\beta\sigma$ with the level-proportional risk price
$\gamma_1\sigma$, plus the analogous product in the size channel, which is
the same structure as the continuous Girsanov computation of the published
paper.
\end{remark}

Two structural readings of \eqref{eq:driftmap} deserve emphasis.
The intercepts $\gamma_0$ and $\eta_0$ move only the linear coefficient, so
the quadratic coefficient is funded exclusively by the slopes of the two
risk prices.
The intercept $d_0$ does not move at all, and the next proposition shows
that this invariance is a property of the recursion class rather than of
the exponential-quadratic family.

\begin{proposition}\label{prop:invariance}
Let the pricing kernel be any positive $\mathcal{F}_{k+1}$-measurable
random variable $M_{k+1}$ with $\E_k[M_{k+1}]=1$ and
$\E_k[M_{k+1}(|z_{k+1}|+|w_{k+1}|)]<\infty$, and suppose the rescaled
adjustment
$\E_k[M_{k+1}(\beta z_{k+1}+\varepsilon w_{k+1})]/\sqrt{\Dt}$ converges
locally uniformly to some function $\Lambda(\sigma)$.
Then the risk-neutral drift satisfies
$\hat b(\sigma)=b(\sigma)+\sigma\Lambda(\sigma)$, and in particular
$\hat b(0)=b(0)=d_0$, equivalently
$\hat\kappa_1\hat\theta=\kappa_1\theta$.
\end{proposition}

\begin{proof}
The change of measure alters the drift of \eqref{eq:sigdyn} by
$\sigma_k\sqrt{\Dt}\,\E_k[M_{k+1}(\beta z_{k+1}+\varepsilon w_{k+1})]$,
which is proportional to $\sigma_k$ for every admissible kernel because the
innovation enters the recursion only through the dispersion term.
Kernels that load on shocks independent of $z$ multiply the adjustment by a
factor with unit expectation and leave the conclusion unchanged.
\end{proof}

\begin{remark}
Proposition~\ref{prop:invariance} is the discrete image of the Girsanov
theorem for the limit diffusion: the drift adjustment is proportional to
the dispersion, which vanishes linearly at zero volatility.
In the published paper the equality $\kappa_1\theta=\hat\kappa_1\hat\theta$
enters as the first of the parameter-matching equations, and the
proposition shows that it is not a convention: no admissible kernel can
break it, at any step size.
\end{remark}

\begin{corollary}\label{cor:theta}
Write the risk-neutral drift as
$\hat b(\sigma)=(\hat\kappa_1+\hat\kappa_2\sigma)(\hat\theta-\sigma)$ with
$\hat\kappa_2>0$.
Then the parameters are determined by the polynomial coefficients through
\begin{equation}\label{eq:thetaQ}
\hat\kappa_2=-\hat d_2,\qquad
\hat\theta
=\frac{\hat d_1+\sqrt{\hat d_1^{\,2}+4\,\hat\kappa_2\,d_0}}
{2\,\hat\kappa_2},
\qquad
\hat\kappa_1=\frac{d_0}{\hat\theta},
\end{equation}
and with $\hat d_1=-(\kappa_1-\kappa_2\theta)$ these are the
risk-transformed parameters of Sepp and Rakhmonov (2023, Section 2).
The level $\hat\theta$ equals the physical level $\theta$ only on the
one-parameter kernel set with $\Lambda_0=(\hat\kappa_2-\kappa_2)\theta$.
\end{corollary}

\begin{proof}
Expanding the factorized form gives $d_0=\hat\kappa_1\hat\theta$,
$\hat d_1=\hat\kappa_2\hat\theta-\hat\kappa_1$ and
$\hat d_2=-\hat\kappa_2$.
Eliminating $\hat\kappa_1$ yields
$\hat\kappa_2\hat\theta^2-\hat d_1\hat\theta-d_0=0$, whose unique positive
root is \eqref{eq:thetaQ} because $\hat\kappa_2>0$ and $d_0>0$.
Setting $\hat\theta=\theta$ in the quadratic and using
$\hat d_1=d_1+\Lambda_0$ with $d_0=\kappa_1\theta$ gives the stated
restriction on $\Lambda_0$.
\end{proof}

Corollary~\ref{cor:theta} settles the parametrization question that
motivated this note.
The triple $(\hat\kappa_1,\hat\kappa_2,\hat\theta)$ obscures the measure
change because all three parameters move together, while the polynomial
coefficients separate cleanly: $\hat d_2$ carries the risk-price slopes,
$\hat d_1$ carries the risk-price levels, and $d_0$ carries nothing.
We therefore parametrize the discrete model, and its reference
implementation, by $(d_0,d_1,d_2)$ with the cross-measure restriction
$d_0^{\Pm}=d_0^{\Qm}$ imposed in any joint estimation from returns and
options.\footnote{Reference implementation planned in the open-source
\texttt{stochvolmodels} package,
\texttt{github.com/ArturSepp/StochVolModels}.}

The transformation formula has one falsifiable consequence worth stating.
On the minimal kernel set of Remark~\ref{rem:esscher}, equation
\eqref{eq:kappa2} reduces to $\hat\kappa_2=\kappa_2+\beta\gamma_1$.
For assets with $\beta>0$, option surfaces must display
$\hat\kappa_2>\kappa_2$ whenever the time-series Sharpe ratio increases in
volatility, before any volatility risk premium is invoked.
For assets with $\beta<0$ the level channel contributes
$-|\beta|\gamma_1$, so an increase of the quadratic rate requires the size
channel with $\eta_1<0$, equivalently $\bar\lambda_1>0$, or a Sharpe ratio
that falls with volatility.
Which channel funds the quadratic drift is an empirical question whose
answer the model links, through \eqref{eq:kappa2}, to the sign of the
volatility beta.

\section{Diffusion limits}\label{sec:limit}

We now state the convergence of the recursion to the continuous model, and
we state it under both measures, because nothing in the argument prefers
one of them.
The physical recursion satisfies the convergence conditions directly, and
the risk-neutral recursion satisfies them after the drift map of
Theorem~\ref{thm:main}.
The conditions are those of Nelson (1990), which descend from the
Markov-chain convergence theory of Stroock and Varadhan, and the
verification is a bookkeeping exercise given the moment expansions above.

\begin{assumption}\label{ass:limit}
The parameters satisfy $\kappa_1,\theta>0$, $\kappa_2\ge0$ and
$\hat\kappa_2\ge0$, the initial values converge, and the recursion is
modified by the localization device of Remark~\ref{rem:trunc}.
\end{assumption}

\begin{theorem}\label{thm:limit}
Under Assumption~\ref{ass:limit} the interpolated processes
$(X^{\Dt},\sigma^{\Dt})$ converge weakly, as $\Dt\to0$, to the unique
solution of the log-normal beta SV model.
Under $\Pm$ the limit is
\begin{equation}\label{eq:limitP}
dX_t=\Big(r+\gamma(\sigma_t)\sigma_t-\tfrac{1}{2}\sigma_t^2\Big)dt
+\sigma_t\,dW^{(0)}_t,
\quad
d\sigma_t=b(\sigma_t)\,dt
+\beta\sigma_t\,dW^{(0)}_t+\varepsilon\sigma_t\,dW^{(1)}_t,
\end{equation}
and under $\Qm$ the limit is the same system with drift
$r-\sigma_t^2/2$ for the log price and
$\hat b(\sigma_t)=(\hat\kappa_1+\hat\kappa_2\sigma_t)
(\hat\theta-\sigma_t)$ for the volatility, with
$W^{(0)}$ and $W^{(1)}$ independent Brownian motions under the respective
measure and parameters given by Corollary~\ref{cor:theta}.
\end{theorem}

\begin{proof}
Under $\Pm$ the per-unit-time conditional means of the increments equal the
drifts of \eqref{eq:limitP} exactly, and under $\Qm$ they converge to the
risk-neutral drifts locally uniformly by Proposition~\ref{prop:qmoments}
and Theorem~\ref{thm:main}.
The conditional covariance matrix of the pair converges to
$\sigma^2\begin{pmatrix}1&\beta\\ \beta&\vartheta^2\end{pmatrix}$ under
both measures, which is the covariance of
$(\sigma dW^{(0)},\,\sigma(\beta dW^{(0)}+\varepsilon dW^{(1)}))$, because
the vector $(z,w)$ is uncorrelated with unit variances and its partial
sums aggregate into two independent Brownian motions.
Fourth conditional moments of the increments are of order $\Dt^2$ times
fourth innovation moments, which are finite for Gaussian innovations, so
they vanish at the required rate on compact sets.
The limit martingale problem is well posed: the volatility equation has a
unique strong solution with values in $(0,\infty)$ despite the
non-Lipschitz drift \cite{SeppRakhmonov2023}, and appending the log-price
equation preserves uniqueness because its coefficients are Lipschitz given
the volatility path.
Convergence then follows from Theorem 2.1 of Nelson (1990).
\end{proof}

\begin{remark}\label{rem:trunc}
Two localization devices make the recursion well defined for all states and
change nothing in the limit.
We truncate the innovation at $\pm\varepsilon_0/\sqrt{\Dt}$ for a small
fixed $\varepsilon_0$, which alters each conditional moment by $o(\Dt)$
under Gaussian tails, and we evaluate the quadratic drift term at
$\sigma\wedge\Dt^{-1/2}$, which is invisible on compact sets.
Weak convergence to a well-posed limit depends on the generator only
through its values on compacts, so both devices are limit-invariant.
The truncation also keeps the volatility strictly positive for small
$\varepsilon_0$, because the drift at zero equals $d_0>0$ while the
dispersion vanishes.
\end{remark}

\begin{remark}\label{rem:equiv}
At every fixed $\Dt$ the one-period kernels define an equivalent measure
change without further conditions.
The equivalence condition of the continuous model, which requires
$\hat\kappa_2\ge0$ in the notation of the published paper, reappears only
through the well-posedness of the limit, so the discrete family
regularizes the measure-equivalence boundary and the constraint binds in
the limit alone.
\end{remark}

\begin{remark}\label{rem:gig}
The stationary density of the limit volatility follows from the scale and
speed measures,
\begin{equation}\label{eq:gig}
p(\sigma)\;\propto\;
\sigma^{\,2\hat d_1/\vartheta^2-2}\,
\exp\!\Big(-\frac{2d_0}{\vartheta^2\,\sigma}
-\frac{2\hat\kappa_2}{\vartheta^2}\,\sigma\Big),
\qquad \sigma>0,
\end{equation}
a generalized inverse Gaussian law, which degenerates to an inverse gamma
law as $\hat\kappa_2\to0$.
This matches the steady state derived by Carr and Willems (2019) for their
quadratic-drift specification and provides a free numerical checkpoint: the
ergodic distribution of the simulated recursion must converge to
\eqref{eq:gig}.
In the linear-drift case $\kappa_2=0$ the tails admit a sharper statement.
At fixed $\Dt$ the recursion is a stochastic recurrence equation, and
Kesten-type results give its stationary law a power tail with index
$\chi_\Dt$ solving
$\E\,|1+d_1\Dt+\sqrt{\Dt}\,(\beta z+\varepsilon w)|^{\chi}=1$.
The limit law is itself power-tailed, an inverse gamma with survival index
$1+2\kappa_1/\vartheta^2$, and $\chi_\Dt$ sits below this value and rises
toward it as the step shrinks.
The companion study confirms the gap and its closure, with
$\chi_\Dt=2.68$, $2.76$ and $2.80$ at weekly, daily and finer steps
against the limit index $2.85$.
The discrete tails are therefore heavier than the limit tails at any
fixed step, by a computable margin rather than by a change of tail class.
\end{remark}

\section{Verification, constraints, and the open problem}\label{sec:open}

Three verification steps preceded this revision, and all three are now on
record in the companion simulation study.
The simulated risk-neutral recursion reproduces the affine-expansion
prices of the continuous model, and the residual at the finest step
measures the expansion itself.
Mean absolute implied-volatility errors are 4 to 5 basis points at the
equity-like parameter set.
At the crypto-like set with $\vartheta=1.80$ they are 23 to 26 basis
points, with a maximum of 58 basis points in the one-month wing.
The full-budget clean-tag rerun reproduces every finest-step strike from the
archived run exactly under the retained common seeds, so the quoted values
are now the clean-tag manuscript numbers.
The kernel algebra is validated end to end.
Prices from the exact risk-neutral law and from likelihood-ratio
reweighting of physical paths agree within Monte Carlo error at every
step size, and the gap of the limit recursion closes at fitted rates of
$0.48$ against the predicted $\sqrt{\Dt}$.
The ergodic density of the simulated volatility matches \eqref{eq:gig}
with relative sup-density errors of $5.63$ and $8.69$ percent at daily
steps, falling to $2.95$ and $6.42$ percent at $\Dt=1/1008$.
The remaining verification item is the cross-measure restriction
$d_0^{\Pm}=d_0^{\Qm}$ on joint return and option data, where it supplies
one overidentifying restriction and one stability diagnostic.

The estimation experiment of Remark~\ref{rem:qmle} quantifies what the
time series alone can deliver, and the answer motivates the
parametrization adopted after Corollary~\ref{cor:theta}.
Across 200 replications at daily steps, the Sharpe slope $\gamma_1$ does
not reach a median absolute $t$-statistic of two even at forty years.
At forty years, $\kappa_2$ carries relative root-mean-square errors of 56
to 100 percent, while the relative errors of the total volatility of
volatility $\vartheta$ and the linearized reversion speed
$\kappa_{\mathrm{lin}}$ are about 4 and 13 percent.
The speed direction aligns with the low-variance eigendirection of the
$(\kappa_1,\kappa_2)$ estimates with cosines above $0.97$; this is a
conditional local ridge diagnostic evaluated at the true $\theta$.
The intercept $d_0$ has a forty-year relative error of 19 to 25 percent,
and the split of the curvature between $\kappa_1$ and $\kappa_2$ is the
ridge.
An oracle ladder quantifies what two restrictions can add.
At forty years, fixing physical $\kappa_2$ and $d_0$ at truth reduces the
$\kappa_1$ RMSE by 79 to 80 percent relative to the unrestricted fit and
by 67 to 69 percent relative to fixing $\kappa_2$ alone; the corresponding
$\theta$ reductions are 10 to 14 percent and 6 to 7 percent.
This is an oracle upper bound on what an option-implied
$\hat\kappa_2$ plus the cross-measure $d_0$ restriction can add to the time
series, not achieved identification from options.

Two constraints bound what any convergence claim may assert.
The $\sqrt{\Dt}$ scaling of the kernel loadings is a modeling choice, and
alternative scalings produce different limits, so every statement above is
relative to the stated scheme \cite{Corradi2000}.
Statistical experiments generated by the recursion and by the diffusion are
not equivalent \cite{Wang2002}, so no claim about transferring inference
between the discrete and the continuous model is made here.

The open problem is the convergence of call prices, and we state it as the
central technical claim of the intended paper.
Weak convergence yields prices of bounded continuous payoffs only.
For $\hat\kappa_2=0$ and $\beta>0$ the discounted limiting stock-price
process is a strict local martingale \cite{Sin1998,LionsMusiela2007}, and
each discrete model prices calls perfectly well at fixed $\Dt$, so the
price family cannot converge to the continuous prices, and the failure
appears only in the limit.
The companion study makes this failure quantitative at a fixed
computational budget.
Every recursion in the family has discounted-spot expectation exactly one
at every step size.
A Monte Carlo estimate with $2^{20}$ paths at one-year maturity
nevertheless falls short of one by about $0.22$ at $\hat\kappa_2=0$ and by
about $0.08$ at $\hat\kappa_2=1/2$, with weak dependence over the finest
two step sizes.
At $\Dt\le1/1008$, estimates for $\hat\kappa_2\ge2$ are statistically
indistinguishable from zero.
The shortfall is not a defect of any fixed-step model but the
finite-budget face of the missing uniform integrability.
Nested budgets from $2^{16}$ to $2^{22}$ paths leave the
$\hat\kappa_2=0$ shortfall near $0.22$ and the largest-budget
$\hat\kappa_2=1/2$ shortfall near $0.09$ at both finest steps.
The latter curve is nonmonotone at the finest step because a rare-path
overshoot at the smallest budget temporarily pushes the empirical mean
above one, while every $\hat\kappa_2=2$ interval contains zero.
This behavior is consistent with a compensating mean contribution lying
in paths missed or irregularly represented at the chosen budget, but an
ordinary bootstrap conditional on observed paths cannot measure unseen
mass.
The conjecture is that a uniform moment bound restores convergence: for
every $\hat\kappa_2\ge\underline{\kappa}>0$ there exists $\varepsilon_1>0$
with
$\sup_{\Dt}\E^{\Qm}\big[(S_T^{\Dt})^{1+\varepsilon_1}\big]<\infty$, which
implies uniform integrability and convergence of call prices.
The corresponding continuous-time result is the true-martingale property
established for the quadratic-drift model.
The discrete conjecture is stronger because it requires a superlinear
moment uniformly in $\Dt$.
In this reading the quadratic rate is the price of admitting a continuum
limit with positive volatility beta: the discrete family with
$\hat\kappa_2=0$ hides the pathology at every fixed step and reveals it
only in the limit.
A second, smaller open item is a discrete moment-generating-function
recursion closed under a quadratic-affine ansatz to the same order as the
affine expansion of the continuous model, which would complete the analogy
with Heston and Nandi (2000).

\begin{thebibliography}{99}

\bibitem[Carr and Willems(2019)]{CarrWillems2019}
Carr, P. and S. Willems (2019).
A lognormal type stochastic volatility model with quadratic drift.
arXiv:1908.07417.

\bibitem[Christoffersen et al.(2013)]{CHJ2013}
Christoffersen, P., S. Heston and K. Jacobs (2013).
Capturing option anomalies with a variance-dependent pricing kernel.
\textit{Review of Financial Studies} 26, 1963--2006.

\bibitem[Corradi(2000)]{Corradi2000}
Corradi, V. (2000).
Reconsidering the continuous time limit of the GARCH(1,1) process.
\textit{Journal of Econometrics} 96, 145--153.
% TODO: verify pages.

\bibitem[Duan(1997)]{Duan1997}
Duan, J.-C. (1997).
Augmented GARCH(p,q) process and its diffusion limit.
\textit{Journal of Econometrics} 79, 97--127.

\bibitem[Heston and Nandi(2000)]{HestonNandi2000}
Heston, S. and S. Nandi (2000).
A closed-form GARCH option valuation model.
\textit{Review of Financial Studies} 13, 585--625.

\bibitem[Lions and Musiela(2007)]{LionsMusiela2007}
Lions, P.-L. and M. Musiela (2007).
Correlations and bounds for stochastic volatility models.
\textit{Annales de l'Institut Henri Poincar\'e (C) Analyse Non Lin\'eaire}
24, 1--16.

\bibitem[Nelson(1990)]{Nelson1990}
Nelson, D. B. (1990).
ARCH models as diffusion approximations.
\textit{Journal of Econometrics} 45, 7--38.

\bibitem[Nelson and Foster(1994)]{NelsonFoster1994}
Nelson, D. B. and D. P. Foster (1994).
Asymptotic filtering theory for univariate ARCH models.
\textit{Econometrica} 62, 1--41.

\bibitem[Sepp and Rakhmonov(2023)]{SeppRakhmonov2023}
Sepp, A. and P. Rakhmonov (2023).
Log-normal stochastic volatility model with quadratic drift.
\textit{International Journal of Theoretical and Applied Finance} 26(8),
2450003.

\bibitem[Sin(1998)]{Sin1998}
Sin, C. A. (1998).
Complications with stochastic volatility models.
\textit{Advances in Applied Probability} 30, 256--268.

\bibitem[Wang(2002)]{Wang2002}
Wang, Y. (2002).
Asymptotic nonequivalence of GARCH models and diffusions.
\textit{Annals of Statistics} 30, 754--783.
% TODO: verify pages.

\bibitem[Zakoian(1994)]{Zakoian1994}
Zakoian, J.-M. (1994).
Threshold heteroskedastic models.
\textit{Journal of Economic Dynamics and Control} 18, 931--955.

\end{thebibliography}

\end{document}
